% Part IV -- Numerical Data, Tables, and Visualization
% Chapter 12: pandas DataFrames and table analysis

\h{pandas DataFrames and Table Analysis}
\label{ch:12}

NumPy arrays organize values by numerical position. Many real tables also need
names for columns and rows. The pandas library provides those labels through a
two-dimensional structure called a \c{DataFrame}. This chapter follows one
small table from creation through selection, cleaning, calculation, and
summary.

\begin{NoteBox}
\textbf{Prerequisites.} You should understand dictionaries, lists, paths, CSV
files, boolean conditions, functions, and NumPy vectorized calculations.
\end{NoteBox}

\begin{tcolorbox}[title=Learning outcomes,colframe=orange,colback=white,arc=2mm]
After completing this chapter, you should be able to:
\begin{itemize}
    \item create and inspect a pandas \c{Series} and \c{DataFrame};
    \item select columns and rows with brackets, \c{loc}, and \c{iloc};
    \item filter rows with one or more boolean conditions;
    \item create calculated columns without writing a row-by-row loop;
    \item detect and handle missing values deliberately;
    \item sort, summarize, and group tabular data; and
    \item read and write a CSV file with a clear index decision.
\end{itemize}
\end{tcolorbox}

\hh{Install once and import when needed}

pandas is an external library. Install it once in the active Python
environment by running this command in a terminal.

\begin{verbatim}
python -m pip install pandas
\end{verbatim}

Use the same Python command that runs your scripts. Replace \texttt{python}
with \texttt{py} or \texttt{python3} if that is the working command on your
computer.

The conventional import name is \c{pd}.

\begin{lstlisting}[
    language=python,
    style=block,
    caption={Import pandas and create a Series},
    label={c12_import_pandas},
    captionpos=b
]
import pandas as pd

values = pd.Series([18.0, 24.0, 31.0])
print(values)
\end{lstlisting}

\hh{Begin with a Series}

A \c{Series} is a one-dimensional sequence with labels. The values appear on
the right and the index labels appear on the left.

\begin{lstlisting}[
    language=python,
    style=block,
    caption={Create a Series with meaningful labels},
    label={c12_labeled_series},
    captionpos=b
]
import pandas as pd

area_by_room = pd.Series(
    [24.0, 18.0, 31.0],
    index=["Lobby", "Office", "Cafe"],
    name="area_m2",
)

print(area_by_room)
print("Cafe area:", area_by_room.loc["Cafe"])
\end{lstlisting}

Labels can communicate meaning more clearly than positions alone. A Series
also keeps a name for the values it stores.

\hh{Create a DataFrame from a dictionary}

A dictionary of equal-length lists is a direct way to build a small table.
Each dictionary key becomes a column name.

\begin{lstlisting}[
    language=python,
    style=block,
    caption={Create a four-row DataFrame},
    label={c12_create_dataframe},
    captionpos=b
]
import pandas as pd

rooms = pd.DataFrame({
    "room": ["Lobby", "Office", "Cafe", "Studio"],
    "zone": ["Public", "Work", "Public", "Work"],
    "area_m2": [24.0, 18.0, 31.0, 26.0],
    "height_m": [4.2, 3.0, 3.6, 3.2],
})

print(rooms)
\end{lstlisting}

Each row is one observation. Each column is one variable measured or recorded
for those observations. Keeping this structure tidy makes later operations
predictable.

\begin{ExplainBox}{A table needs both a grain and a schema}
The \emph{grain} states what one row represents: one room, one sensor reading,
one building, or one event. The \emph{schema} states what each column means,
which type it should contain, its units, and whether missing values are valid.
Without an explicit grain, duplicate-looking rows cannot be interpreted.
Without a schema, a column named \c{value} cannot be validated or compared.
\end{ExplainBox}

For the \c{rooms} table, the grain is one row per room. The schema says that
\c{room} and \c{zone} are categorical text, while \c{area_m2} and
\c{height_m} are measured numerical values with stated units. If
\c{area_m2} loads as text because one cell contains \c{"unknown"}, vectorized
arithmetic is no longer trustworthy. Inspection is therefore evidence that
the loaded data matches the analytical claim, not merely a preliminary coding
ritual.

\bookfigure[0.94\textwidth]{Figures/Generated/ch12_dataframe_pipeline.pdf}
{A readable pandas analysis makes loading, inspection, cleaning, transformation, and communication visible.}
{fig:ch12-dataframe-pipeline}

\hh{Inspect before analyzing}

Do not assume that a loaded table has the expected size or data types. Begin
with a few focused checks.

\begin{lstlisting}[
    language=python,
    style=block,
    caption={Inspect structure and a small sample},
    label={c12_inspect_dataframe},
    captionpos=b
]
import pandas as pd

rooms = pd.DataFrame({
    "room": ["Lobby", "Office", "Cafe", "Studio"],
    "zone": ["Public", "Work", "Public", "Work"],
    "area_m2": [24.0, 18.0, 31.0, 26.0],
    "height_m": [4.2, 3.0, 3.6, 3.2],
})

print("Shape:", rooms.shape)
print("Columns:", rooms.columns.tolist())
print("Data types:\n", rooms.dtypes)
print("First rows:\n", rooms.head(2))
\end{lstlisting}

The shape reports rows first and columns second. \c{head()} displays a sample;
it does not prove that the rest of a large table is clean.

\hh{Select one or several columns}

Single brackets with one column name return a Series. Double brackets with a
list of names return a DataFrame.

\begin{lstlisting}[
    language=python,
    style=block,
    caption={Distinguish Series and DataFrame column selections},
    label={c12_select_columns},
    captionpos=b
]
import pandas as pd

rooms = pd.DataFrame({
    "room": ["Lobby", "Office", "Cafe"],
    "area_m2": [24.0, 18.0, 31.0],
    "height_m": [4.2, 3.0, 3.6],
})

areas = rooms["area_m2"]
labels_and_areas = rooms[["room", "area_m2"]]

print(type(areas).__name__)
print(labels_and_areas)
\end{lstlisting}

Choose the form based on the result you need. Do not add extra brackets only
because both versions happen to print like tables.

\hh{Select rows by label or position}

Use \c{loc} for labels and \c{iloc} for integer positions. Making that choice
explicit prevents confusion when an index is not a simple sequence.

\begin{lstlisting}[
    language=python,
    style=block,
    caption={Select rows and cells with loc and iloc},
    label={c12_loc_and_iloc},
    captionpos=b
]
import pandas as pd

rooms = pd.DataFrame(
    {
        "area_m2": [24.0, 18.0, 31.0],
        "height_m": [4.2, 3.0, 3.6],
    },
    index=["Lobby", "Office", "Cafe"],
)

print("By label:\n", rooms.loc["Office"])
print("By position:\n", rooms.iloc[1])
print("Selected cells:\n", rooms.loc[["Lobby", "Cafe"], ["area_m2"]])
\end{lstlisting}

Unlike an ordinary Python slice, a label slice used with \c{loc} includes its
ending label. Position slices used with \c{iloc} stop before the ending
position.

\begin{SyntaxBox}{Use loc when the selection has meaning}
The pattern \c{df.loc[row_selection, column_selection]} states both
dimensions. For example, \c{rooms.loc[rooms["area_m2"] >= 25, ["height_m",
"area_m2"]]} keeps matching rows and only the named columns. Use \c{iloc}
when the selection is intentionally based on integer position.
\end{SyntaxBox}

\begin{NoteBox}
\textbf{Stop and predict.} If the index labels are \c{"A"}, \c{"B"}, and
\c{"C"}, which rows are returned by \c{df.loc["A":"B"]}? Which positions are
returned by \c{df.iloc[0:2]}?
\end{NoteBox}

\hh{Filter rows with conditions}

Build a boolean Series, also called a boolean mask, then place it inside
brackets to keep matching rows.

\begin{lstlisting}[
    language=python,
    style=block,
    caption={Filter a DataFrame with one condition},
    label={c12_filter_rows},
    captionpos=b
]
import pandas as pd

rooms = pd.DataFrame({
    "room": ["Lobby", "Office", "Cafe", "Studio"],
    "zone": ["Public", "Work", "Public", "Work"],
    "area_m2": [24.0, 18.0, 31.0, 26.0],
})

large_rooms = rooms[rooms["area_m2"] >= 25.0]
print(large_rooms)
\end{lstlisting}

Combine conditions with \c{&} or \c{|}, and put each comparison in
parentheses.

\begin{lstlisting}[
    language=python,
    style=block,
    caption={Filter with two conditions},
    label={c12_filter_multiple_conditions},
    captionpos=b
]
import pandas as pd

rooms = pd.DataFrame({
    "room": ["Lobby", "Office", "Cafe", "Studio"],
    "zone": ["Public", "Work", "Public", "Work"],
    "area_m2": [24.0, 18.0, 31.0, 26.0],
})

selected = rooms[
    (rooms["zone"] == "Public") & (rooms["area_m2"] >= 25.0)
]
print(selected)
\end{lstlisting}

\hh{Create a calculated column}

Column arithmetic is vectorized. The calculation aligns values by row.

\begin{lstlisting}[
    language=python,
    style=block,
    caption={Add a volume column from two existing columns},
    label={c12_calculated_column},
    captionpos=b
]
import pandas as pd

rooms = pd.DataFrame({
    "room": ["Lobby", "Office", "Cafe"],
    "area_m2": [24.0, 18.0, 31.0],
    "height_m": [4.2, 3.0, 3.6],
})

rooms["volume_m3"] = rooms["area_m2"] * rooms["height_m"]
print(rooms)
\end{lstlisting}

The new column name includes the unit. That small naming choice prevents
mistakes when several measurement systems appear in one table.

\hh{Update selected rows with loc}

Use \c{loc[row condition, column name]} when changing only part of a column.

\begin{lstlisting}[
    language=python,
    style=block,
    caption={Assign values to selected rows explicitly},
    label={c12_assign_with_loc},
    captionpos=b
]
import pandas as pd

rooms = pd.DataFrame({
    "room": ["Lobby", "Office", "Cafe"],
    "zone": ["Public", "Work", "Public"],
    "status": ["Review", "Review", "Review"],
})

public_mask = rooms["zone"] == "Public"
rooms.loc[public_mask, "status"] = "Checked"

print(rooms)
\end{lstlisting}

Avoid changing a temporary filtered object with chained brackets. One \c{loc}
operation clearly states both the rows and the column to update.

\hh{Find missing values}

Missing values require a decision. First count them; then decide whether to
fill, remove, or investigate them.

\begin{lstlisting}[
    language=python,
    style=block,
    caption={Detect and count missing values},
    label={c12_detect_missing_values},
    captionpos=b
]
import pandas as pd

rooms = pd.DataFrame({
    "room": ["Lobby", "Office", "Cafe", "Studio"],
    "area_m2": [24.0, None, 31.0, None],
    "zone": ["Public", "Work", "Public", "Work"],
})

print(rooms.isna())
print("Missing by column:\n", rooms.isna().sum())
\end{lstlisting}

Do not test missing values with \c{== None} or \c{== np.nan}. Use \c{isna()}
or \c{notna()} so pandas handles its missing-value representations correctly.

\hh{Handle missing values deliberately}

Filling and removing answer different questions. Filling preserves the row but
introduces an assumed value. Removing avoids that assumption but loses data.

\begin{lstlisting}[
    language=python,
    style=block,
    caption={Compare filling with dropping incomplete rows},
    label={c12_fill_missing_values},
    captionpos=b
]
import pandas as pd

rooms = pd.DataFrame({
    "room": ["Lobby", "Office", "Cafe", "Studio"],
    "area_m2": [24.0, None, 31.0, None],
})

median_area = rooms["area_m2"].median()
filled = rooms.copy()
filled["area_m2"] = filled["area_m2"].fillna(median_area)

complete = rooms.dropna(subset=["area_m2"])

print("Filled:\n", filled)
print("Complete rows:\n", complete)
\end{lstlisting}

Keep the original DataFrame unchanged while comparing cleaning choices. Record
the rule used to fill values so a reader can evaluate the assumption.

\hh{Sort without losing the original order}

\c{sort_values()} returns a sorted DataFrame unless you explicitly request an
in-place change. Assign the result to a descriptive name.

\begin{lstlisting}[
    language=python,
    style=block,
    caption={Sort rows by a numerical column},
    label={c12_sort_values},
    captionpos=b
]
import pandas as pd

rooms = pd.DataFrame({
    "room": ["Lobby", "Office", "Cafe", "Studio"],
    "area_m2": [24.0, 18.0, 31.0, 26.0],
})

largest_first = rooms.sort_values("area_m2", ascending=False)
print(largest_first)
\end{lstlisting}

The index stays attached to its row during sorting. Use \c{reset_index()} only
when a new sequential index is actually useful.

\hh{Calculate column summaries}

Series methods calculate focused statistics. \c{describe()} gives a compact
overview of numerical columns.

\begin{lstlisting}[
    language=python,
    style=block,
    caption={Summarize numerical columns},
    label={c12_summary_statistics},
    captionpos=b
]
import pandas as pd

rooms = pd.DataFrame({
    "area_m2": [24.0, 18.0, 31.0, 26.0],
    "height_m": [4.2, 3.0, 3.6, 3.2],
})

print("Total area:", rooms["area_m2"].sum())
print("Mean area:", rooms["area_m2"].mean())
print(rooms.describe())
\end{lstlisting}

Summary statistics support inspection; they do not explain why a value is
unusual. Return to the relevant rows when a minimum or maximum needs context.

\hh{Group rows and aggregate}

Grouping follows three steps: split rows into categories, apply a calculation,
and combine the results.

\begin{lstlisting}[
    language=python,
    style=block,
    caption={Count and summarize areas by category},
    label={c12_groupby_summary},
    captionpos=b
]
import pandas as pd

rooms = pd.DataFrame({
    "room": ["Lobby", "Office", "Cafe", "Studio", "Meeting"],
    "zone": ["Public", "Work", "Public", "Work", "Work"],
    "area_m2": [24.0, 18.0, 31.0, 26.0, 20.0],
})

area_by_zone = rooms.groupby("zone", as_index=False).agg(
    room_count=("room", "count"),
    total_area_m2=("area_m2", "sum"),
    mean_area_m2=("area_m2", "mean"),
)

print(area_by_zone)
\end{lstlisting}

Read the expression from left to right: group by \c{zone}, then calculate three
named summaries. Each pair states an output column name, a source column, and
an aggregation operation.

\bookfigure[0.94\textwidth]{Figures/Generated/ch12_dataframe_before_after.pdf}
{A raw room table becomes an inspectable floor-level summary after cleaning, grouping, and aggregation.}
{fig:ch12-before-after}

Always compare the input and result of a transformation. Row counts, missing
values, grouping keys, and totals provide simple invariants that can reveal a
mistake before it reaches a chart or exported file.

\hh{Read and write CSV data}

Use \c{to_csv()} to save a DataFrame and \c{read_csv()} to restore it. Decide
whether the DataFrame index is data that belongs in the file.

\begin{lstlisting}[
    language=python,
    style=block,
    caption={Write and read a small CSV file},
    label={c12_csv_round_trip},
    captionpos=b
]
from pathlib import Path
import pandas as pd

folder = Path("study_data")
folder.mkdir(exist_ok=True)
csv_path = folder / "rooms_pandas.csv"

rooms = pd.DataFrame({
    "room": ["Lobby", "Office", "Cafe"],
    "area_m2": [24.0, 18.0, 31.0],
})

rooms.to_csv(csv_path, index=False)
restored = pd.read_csv(csv_path)

print(restored)
print(restored.dtypes)
\end{lstlisting}

Here the default row index is only an internal counter, so \c{index=False}
keeps it out of the CSV. A meaningful index may deserve its own named column.

\hh{Build a readable analysis sequence}

Keep inspection, cleaning, calculation, and summary as visible stages. Separate
names let you compare stages and locate a mistake.

\begin{lstlisting}[
    language=python,
    style=block,
    caption={Move through clear analysis stages},
    label={c12_analysis_sequence},
    captionpos=b
]
import pandas as pd

raw = pd.DataFrame({
    "room": ["Lobby", "Office", "Cafe", "Studio", "Meeting"],
    "zone": ["Public", "Work", "Public", "Work", "Work"],
    "area_m2": [24.0, None, 31.0, 26.0, 20.0],
    "height_m": [4.2, 3.0, 3.6, 3.2, 3.0],
})

clean = raw.dropna(subset=["area_m2"]).copy()
clean["volume_m3"] = clean["area_m2"] * clean["height_m"]
summary = clean.groupby("zone", as_index=False).agg(
    room_count=("room", "count"),
    total_area_m2=("area_m2", "sum"),
    mean_volume_m3=("volume_m3", "mean"),
)

print(summary)
\end{lstlisting}

Named aggregation gives the result descriptive column names. The cleaning rule
is also visible: only rows missing \c{area_m2} are removed.

\begin{TryBox}{Narrate an analysis pipeline}
Run the staged example one statement at a time. After each stage, print
\c{shape} and \c{head()}. In a short comment, state what rows or columns could
change at that stage. This turns a compact method chain into an auditable
sequence you can debug and explain.
\end{TryBox}

\hh{Common mistakes}

\begin{itemize}
    \item \textbf{Skipping inspection}: check shape, columns, data types, and a
    small sample before analysis.
    \item \textbf{Confusing a Series with a DataFrame}: one column name returns
    a Series; a list of column names returns a DataFrame.
    \item \textbf{Mixing \c{loc} and \c{iloc}}: use labels with \c{loc} and
    positions with \c{iloc}.
    \item \textbf{Combining conditions without parentheses}: parenthesize each
    comparison before \c{&} or \c{|}.
    \item \textbf{Using chained assignment}: select rows and the destination
    column together with \c{loc}.
    \item \textbf{Filling missing data automatically}: state and examine the
    assumption behind any replacement value.
    \item \textbf{Saving an unwanted index}: use \c{index=False} when the index
    is not part of the data.
\end{itemize}

\hh{Practice ladder}

\hhh{Level 0 -- Read and predict}

Predict the type names and the number of rows in \c{selected}.

\begin{lstlisting}[
    language=python,
    style=block,
    caption={Predict selection types and size},
    label={c12_practice_level0},
    captionpos=b
]
import pandas as pd

data = pd.DataFrame({
    "name": ["A", "B", "C"],
    "value": [10, 25, 18],
})
one_column = data["value"]
selected = data[data["value"] >= 18]

print(type(one_column).__name__)
print(type(selected).__name__)
print(len(selected))
\end{lstlisting}

\hhh{Level 1 -- Modify}

Change the filter to keep rows from the \c{"South"} zone with an area of at
least \c{20.0} square meters.

\begin{lstlisting}[
    language=python,
    style=block,
    caption={Modify a row filter},
    label={c12_practice_level1},
    captionpos=b
]
import pandas as pd

rooms = pd.DataFrame({
    "room": ["A", "B", "C", "D"],
    "zone": ["North", "South", "South", "North"],
    "area_m2": [18.0, 24.0, 19.0, 28.0],
})

selected = rooms[rooms["area_m2"] >= 18.0]
print(selected)
\end{lstlisting}

\hhh{Level 2 -- Complete}

Create a DataFrame with columns \c{item}, \c{quantity}, and \c{unit_cost}.
Add a \c{total_cost} column, then print only \c{item} and \c{total_cost}.

\hhh{Level 3 -- Apply}

Create a five-row table containing \c{room}, \c{floor}, and \c{area_m2}. Group
by floor and report the room count and total area for each floor.

\textbf{Hints}
\begin{itemize}
    \item \textbf{Hint 1}: Use one integer floor value per row.
    \item \textbf{Hint 2}: Group the rows by floor without keeping the group as
    the index.
    \item \textbf{Hint 3}: Named aggregation can create clear output column names.
\end{itemize}

\hhh{Level 4 -- Challenge}

Write \c{clean_area_table(data)}. It should return a copy that removes rows
with missing \c{area_m2}, keeps only positive areas, and adds
\c{area_ft2 = area_m2 * 10.7639}. Do not change the input DataFrame.

\hh{Practice solutions}

\textbf{Level 0:} The types are \c{Series} and \c{DataFrame}; two rows remain.

\textbf{Level 2}

\begin{lstlisting}[
    language=python,
    style=block,
    caption={One solution to Level 2},
    label={c12_solution_level2},
    captionpos=b
]
import pandas as pd

costs = pd.DataFrame({
    "item": ["Panel", "Fastener", "Sealant"],
    "quantity": [12, 80, 6],
    "unit_cost": [45.0, 0.75, 18.0],
})
costs["total_cost"] = costs["quantity"] * costs["unit_cost"]

print(costs[["item", "total_cost"]])
\end{lstlisting}

\textbf{Level 3}

\begin{lstlisting}[
    language=python,
    style=block,
    caption={One solution to Level 3},
    label={c12_solution_level3},
    captionpos=b
]
import pandas as pd

rooms = pd.DataFrame({
    "room": ["Lobby", "Office", "Cafe", "Studio", "Meeting"],
    "floor": [1, 1, 1, 2, 2],
    "area_m2": [24.0, 18.0, 31.0, 26.0, 20.0],
})

summary = rooms.groupby("floor", as_index=False).agg(
    room_count=("room", "count"),
    total_area_m2=("area_m2", "sum"),
)
print(summary)
\end{lstlisting}

\textbf{Level 4}

\begin{lstlisting}[
    language=python,
    style=block,
    caption={One solution to Level 4},
    label={c12_solution_level4},
    captionpos=b
]
import pandas as pd

def clean_area_table(data):
    clean = data.dropna(subset=["area_m2"]).copy()
    clean = clean[clean["area_m2"] > 0].copy()
    clean["area_ft2"] = clean["area_m2"] * 10.7639
    return clean

raw = pd.DataFrame({
    "room": ["A", "B", "C", "D"],
    "area_m2": [18.0, None, -4.0, 25.0],
})

cleaned = clean_area_table(raw)
print(cleaned)
print("Original columns:", raw.columns.tolist())
\end{lstlisting}

\hh{Chapter checkpoint}

\begin{enumerate}
    \item What does a DataFrame row usually represent?
    \item What is the result of selecting one column with one name?
    \item When should you use \c{loc} rather than \c{iloc}?
    \item Why must combined pandas conditions be parenthesized?
    \item How can you assign to selected rows and one column clearly?
    \item Which method detects missing values?
    \item Why is filling a missing value a substantive decision?
    \item What three ideas describe a grouped aggregation?
    \item Why might \c{to_csv()} use \c{index=False}?
    \item Why keep raw, clean, and summary tables in separate variables?
\end{enumerate}

\hh{Checkpoint answers}

\begin{enumerate}
    \item One observation, such as one room or one measurement.
    \item A Series.
    \item Use \c{loc} when selecting by labels; use \c{iloc} for positions.
    \item Parentheses define each comparison before element-wise \c{&} or \c{|} joins them.
    \item Use \c{df.loc[row_condition, "column"] = value}.
    \item \c{isna()}.
    \item The replacement introduces an assumption about unknown data.
    \item Split rows into groups, apply a calculation, and combine the results.
    \item To avoid writing an index that is only an internal row counter.
    \item The stages remain inspectable, and the source data is preserved.
\end{enumerate}

\begin{tcolorbox}[title=Chapter summary,colframe=orange,colback=white,arc=2mm]
pandas adds row and column labels to tabular data. Inspect a DataFrame before
analysis, select rows and columns explicitly, handle missing values according
to a stated rule, and keep analysis stages visible. The resulting summaries
are ready to communicate through well-chosen charts.
\end{tcolorbox}

